% META: {"id": "TSPE-TRIGO-CO-021", "chapitre": "TSPE-TRIGONOMETRIE", "type_objet": "corrige", "exercice_ref": "TSPE-TRIGO-EX-021", "capacites_codes": ["C2"], "sources_inspiration": [], "status": "generated"}
% BEGIN-VERIFY
% from sympy import symbols, sin, cos, diff, pi, sqrt, solveset, Interval, simplify
% x = symbols('x', real=True)
% f = sin(x) + cos(x)
% assert simplify(diff(f, x) - (cos(x) - sin(x))) == 0
% zeros = solveset(cos(x) - sin(x), x, Interval(0, 2 * pi))
% assert zeros == {pi / 4, 5 * pi / 4}
% assert simplify(f.subs(x, pi / 4) - sqrt(2)) == 0
% assert simplify(f.subs(x, 5 * pi / 4) + sqrt(2)) == 0
% assert simplify(f - sqrt(2) * sin(x + pi / 4)) == 0
% assert diff(f, x).subs(x, 0) > 0
% END-VERIFY
\begin{corrige}{TSPE-TRIGO-EX-021}
\textbf{1.} $f'(x)=\cos x-\sin x$, qui s'annule lorsque $\tan x=1$, soit
$x=\tfrac{\pi}{4}$ ou $x=\tfrac{5\pi}{4}$ sur $[0\,;\,2\pi]$.

\textbf{2.} $f'$ est positive sur $\left[0\,;\,\tfrac{\pi}{4}\right]$,
négative sur $\left[\tfrac{\pi}{4}\,;\,\tfrac{5\pi}{4}\right]$, positive sur
$\left[\tfrac{5\pi}{4}\,;\,2\pi\right]$. $f$ croît, décroît, puis croît.

\textbf{3.} Le maximum est atteint en $x=\tfrac{\pi}{4}$ et vaut
$f\left(\tfrac{\pi}{4}\right)=\tfrac{\sqrt2}{2}+\tfrac{\sqrt2}{2}=\sqrt2$.
On peut le retrouver en écrivant
$f(x)=\sqrt2\sin\left(x+\tfrac{\pi}{4}\right)$.
\end{corrige}
